\documentclass[a4paper,11pt]{article}
\usepackage[utf8]{inputenc}
\usepackage[french]{babel}
\usepackage[T1]{fontenc}
\usepackage{amsmath,amssymb,amsthm}
\usepackage{geometry}
\geometry{margin=2cm}
\usepackage{enumitem}
\usepackage{hyperref}
\usepackage{booktabs}
\usepackage{array}
\usepackage{xcolor}
\usepackage{listings}
\usepackage{titlesec}
\usepackage{fancyhdr}
\usepackage{lastpage}
\usepackage{graphicx}
\usepackage{etoolbox}
\AtBeginDocument{\shorthandoff{;:!?}}

\titleformat{\section}{\large\bfseries\color{blue!60!black}}{}{0pt}{}[\vspace{-0.5ex}\rule{\textwidth}{0.4pt}]
\titleformat{\subsection}{\bfseries\color{blue!40!black}}{}{0pt}{}
\titleformat{\subsubsection}{\itshape\color{blue!30!black}}{}{0pt}{}

\lstdefinestyle{monstyle}{
  frame=single,
  numbers=left,
  numbersep=5pt,
  breaklines=true,
  basicstyle=\small\ttfamily,
  backgroundcolor=\color{gray!5},
  rulecolor=\color{gray!40},
  columns=fullflexible,
  keepspaces=true,
  showstringspaces=false
}
\lstset{style=monstyle}

\pagestyle{fancy}
\fancyhf{}
\fancyhead[L]{\small STA211 -- M\'ethodes non-supervis\'ees}
\fancyhead[R]{\small Fiche r\'ecapitulative}
\fancyfoot[C]{\thepage/\pageref{LastPage}}
\renewcommand{\headrulewidth}{0.4pt}

\theoremstyle{definition}
\newtheorem*{definition}{D\'efinition}
\theoremstyle{remark}
\newtheorem*{remark}{Remarque}
\newtheorem*{important}{\`A retenir}

\title{\textbf{Fiche r\'ecapitulative -- STA211}\\[4pt]
\large M\'ethodes non-supervis\'ees (introduction)}
\author{Vincent Audigier, Nd\`eye Niang}
\date{12 juin 2026}

\begin{document}
\maketitle
\thispagestyle{empty}

\begin{abstract}
Cette fiche pr\'esente une introduction aux m\'ethodes non-supervis\'ees en
fouille de donn\'ees. Elle couvre les m\'ethodes de classification
(partitionnement et hi\'erarchique), la classification de variables, les
r\`egles d'association, et les m\'ethodes d'analyse factorielle (ACP, ACM,
AFDM, m\'ethodes multi-blocs). Les cartes de Kohonen font l'objet d'une
fiche d\'edi\'ee.
\end{abstract}

\vfill
\tableofcontents
\newpage

% ===================================================================
\section{Contexte et panorama}
% ===================================================================

Les m\'ethodes non-supervis\'ees se distinguent des m\'ethodes
supervis\'ees par l'absence de variable cible. L'objectif est de
d\'ecouvrir des structures, des groupes, ou des associations au sein
des donn\'ees, sans connaissance a priori.

On distingue trois grandes familles~:

\begin{enumerate}
  \item \textbf{Classification} (ou clustering)~: identifier des groupes
    homog\`enes d'individus ou de variables.
  \item \textbf{Recherche de r\`egles d'association}~: identifier des
    associations fr\'equentes entre items (ex.~paniers d'achat).
  \item \textbf{Analyse factorielle}~: r\'esumer les donn\'ees par des
    graphiques de faible dimension (2D/3D) pour visualiser ressemblances
    entre individus et relations entre variables.
\end{enumerate}

Les m\'ethodes de classification peuvent \^etre bas\'ees sur des
distances (approches g\'eom\'etriques) ou sur des mod\`eles.
Les m\'ethodes d'analyse factorielle sont compl\'ementaires~: elles
peuvent servir de pr\'e-traitement (transformation de variables
qualitatives en quantitatives, r\'eduction de dimension, am\'elioration
de la partition via le \emph{tandem clustering}).

\begin{remark}
  Les cartes de Kohonen (SOM), les m\'ethodes de classification
  d'individus classiques (k-means, CAH) et les analyses factorielles
  de base (ACP, ACM, AFDM) sont suppos\'ees connues. Cette fiche
  pr\'esente une synth\`ese et les prolongements (classification de
  variables, m\'ethodes multi-blocs).
\end{remark}

% ===================================================================
\section{Classification des individus}
% ===================================================================

\subsection{M\'ethodes de partitionnement}

Les m\'ethodes de partitionnement consistent \`a d\'efinir une partition
de l'ensemble des individus en $K$ groupes (d\'efini \`a l'avance)
telle que la dispersion intra-groupe soit minimale et la dispersion
inter-groupe maximale.

\subsubsection{K-means}

L'algorithme des $K$-means (MacQueen, 1967; Forgy, 1965) minimise
l'inertie intra-classe~:

\begin{equation}
W(C) = \sum_{k=1}^{K} \sum_{i \in C_k} \| x_i - \mu_k \|^2
\end{equation}

o\`u $\mu_k$ est le centro\"ide du groupe $C_k$.

Algorithme~:
\begin{enumerate}
  \item Choisir $K$ centres initiaux (al\'eatoirement ou via
    $K$-means++).
  \item Affecter chaque individu au centre le plus proche (distance
    euclidienne).
  \item Recalculer les centres comme moyenne des individus du groupe.
  \item R\'ep\'eter 2--3 jusqu'\`a convergence (stabilit\'e des
    affectations).
\end{enumerate}

\begin{important}
  $K$-means est sensible aux centres initiaux et aux valeurs
  atypiques. On utilise $K$-means++ pour une initialisation robuste.
\end{important}

\subsubsection{Partitionnement autour des m\'edo\"ides (PAM)}

Alternative robuste \`a $K$-means utilisant des m\'edo\"ides
(individus r\'eels plut\^ot que des moyennes) comme centres.
Accepte toute matrice de dissimilarit\'e.

\subsection{M\'ethodes hi\'erarchiques}

Les m\'ethodes hi\'erarchiques construisent une hi\'erarchie de
partitions, repr\'esent\'ee par un dendrogramme. L'utilisateur
choisit la partition finale en coupant l'arbre \`a un niveau
donn\'e.

\subsubsection{CAH (Classification Ascendante Hi\'erarchique)}

Algorithme agglom\'eratif~:
\begin{enumerate}
  \item Initialiser chaque individu comme un groupe.
  \item Fusionner les deux groupes les plus proches (selon un crit\`ere
    de saut minimal).
  \item R\'ep\'eter jusqu'\`a obtenir un seul groupe.
\end{enumerate}

Principaux crit\`eres de saut (liens)~:

\begin{tabular}{ll}
\toprule
Crit\`ere & Distance entre groupes $A$ et $B$ \\
\midrule
Saut minimum & $\displaystyle \min_{i \in A, j \in B} d(i,j)$ \\
Saut maximum & $\displaystyle \max_{i \in A, j \in B} d(i,j)$ \\
Lien moyen & $\displaystyle \frac{1}{|A||B|} \sum_{i \in A} \sum_{j \in B} d(i,j)$ \\
Ward & $\displaystyle \frac{|A||B|}{|A|+|B|} \| \mu_A - \mu_B \|^2$ \\
\bottomrule
\end{tabular}

\begin{remark}
  Le crit\`ere de Ward est le plus utilis\'e~: il minimise la
  perte d'inertie intra-classe \`a chaque fusion.
\end{remark}

\subsubsection{M\'ethode divisive}

Approche inverse de la CAH~: on part d'un seul groupe que l'on
divise r\'ecursivement. Moins utilis\'ee car plus co\^uteuse
($O(2^n)$ sans heuristiques).

\subsection{Choix du nombre de groupes}

Plusieurs crit\`eres pour d\'eterminer $K$~:
\begin{itemize}
  \item \textbf{M\'ethode du coude} (elbow)~: inertie intra-classe
    en fonction de $K$; on cherche le point d'inflexion.
  \item \textbf{Indice de silhouette}~: mesure la coh\'esion interne
    et la s\'eparation des groupes $\in [-1, 1]$.
  \item \textbf{Indice de Davies--Bouldin}~: rapport de la dispersion
    intra-groupe \`a la distance inter-groupes (minimiser).
  \item \textbf{Saut de variance} (\emph{gap statistic})~: compare
    l'inertie observ\'ee \`a celle attendue sous une distribution
    de r\'ef\'erence.
\end{itemize}

% ===================================================================
\section{Classification de variables}
% ===================================================================

La classification de variables est similaire dans son principe \`a
celle des individus, mais pr\'esente des particularit\'es li\'ees \`a
la nature des variables (quantitatives, qualitatives, mixtes).

\subsection{Mesures de similarit\'e entre variables}

\begin{itemize}
  \item \textbf{Variables quantitatives}~: corr\'elation de Pearson,
    corr\'elation de Spearman, information mutuelle.
  \item \textbf{Variables qualitatives}~: V de Cramer, $\phi$ de
    contingence, information mutuelle normalis\'ee.
  \item \textbf{Variables mixtes}~: coefficient de corr\'elation
    de Pearson entre les rapports de corr\'elation (via l'AFDM).
\end{itemize}

\subsection{Application}

On peut appliquer une CAH ou $K$-means sur la matrice de similarit\'e
entre variables apr\`es transformation en dissimilarit\'e
($d = 1 - s$). L'objectif est de former des groupes de variables
redondantes pour les r\'esumer.

\begin{important}
  La classification de variables est utile pour la s\'election de
  variables (garder une variable repr\'esentative par groupe) et
  pour l'interpr\'etation des r\'esultats d'analyses factorielles.
\end{important}

% ===================================================================
\section{Recherche de r\`egles d'association}
% ===================================================================

Les r\`egles d'association sont tr\`es utilis\'ees en marketing
(analyse de paniers d'achat). Elles consistent \`a identifier des
associations fr\'equentes entre items.

\subsection{D\'efinitions}

Soit un ensemble d'items $\mathcal{I} = \{i_1, \dots, i_m\}$ et un
ensemble de transactions $\mathcal{T} = \{t_1, \dots, t_n\}$ o\`u
chaque transaction $t \subseteq \mathcal{I}$.

Une r\`egle d'association est une implication $X \Rightarrow Y$ o\`u
$X, Y \subseteq \mathcal{I}$ et $X \cap Y = \emptyset$.

\begin{itemize}
  \item \textbf{Support}~: fr\'equence de l'itemset $X$~:
    $\displaystyle \text{supp}(X) = \frac{|\{t \in \mathcal{T} : X \subseteq t\}|}{n}$
  \item \textbf{Confiance}~: probabilit\'e conditionnelle de $Y$
    sachant $X$~:
    $\displaystyle \text{conf}(X \Rightarrow Y) = \frac{\text{supp}(X \cup Y)}{\text{supp}(X)}$
  \item \textbf{Lift}~: rapport entre la confiance et le support de $Y$~:
    $\displaystyle \text{lift}(X \Rightarrow Y) = \frac{\text{supp}(X \cup Y)}{\text{supp}(X)\,\text{supp}(Y)}$
\end{itemize}

\begin{remark}
  Un lift $> 1$ indique une association positive (les items
  apparaissent plus souvent ensemble que sous l'ind\'ependance).
  Un lift $< 1$ indique une association n\'egative.
\end{remark}

\subsection{Algorithme Apriori}

L'algorithme Apriori (Agrawal \& Srikant, 1994) proc\`ede en deux
\'etapes~:

\begin{enumerate}
  \item \textbf{G\'en\'eration des itemsets fr\'equents}~: on
    parcourt les itemsets par taille croissante en \'elaguant ceux
    dont le support est inf\'erieur \`a un seuil $\text{supp}_{\min}$
    (propri\'et\'e d'anti-monotonie~: tout sur-ensemble d'un itemset
    non fr\'equent est non fr\'equent).
  \item \textbf{G\'en\'eration des r\`egles}~: \`a partir des
    itemsets fr\'equents, on g\'en\`ere les r\`egles $X \Rightarrow Y$
    v\'erifiant $\text{conf} \geq \text{conf}_{\min}$.
\end{enumerate}

\subsection{Qualit\'e des r\`egles}

Plusieurs mesures pour \'evaluer une r\`egle~:
\begin{itemize}
  \item Confiance (d\'efaut)
  \item Lift
  \item Conviction~: $\displaystyle \frac{1 - \text{supp}(Y)}{1 - \text{conf}(X \Rightarrow Y)}$
  \item Levier~: $\text{supp}(X \cup Y) - \text{supp}(X)\,\text{supp}(Y)$
\end{itemize}

% ===================================================================
\section{Analyse factorielle}
% ===================================================================

Les m\'ethodes d'analyse factorielle sont des approches g\'eom\'etriques
qui r\'esument les donn\'ees sous forme de graphiques de dimension
r\'eduite (g\'en\'eralement 2 ou 3).

\subsection{M\'ethodes classiques}

\begin{itemize}
  \item \textbf{ACP (Analyse en Composantes Principales)}~: pour
    des variables quantitatives. R\'eduit la dimension en projetant
    les donn\'ees sur des axes maximisant la variance (inertie).
  \item \textbf{ACM (Analyse des Correspondances Multiples)}~: pour
    des variables qualitatives. \'Etend l'analyse des correspondances
    \`a plus de deux variables.
  \item \textbf{AFDM (Analyse Factorielle de Donn\'ees Mixtes)}~:
    pour des donn\'ees mixtes (quantitatives + qualitatives).
    Combine ACP et ACM en une analyse unique.
\end{itemize}

\subsection{Utilit\'e en pr\'e-traitement}

L'analyse factorielle est compl\'ementaire aux m\'ethodes de
classification~:

\begin{itemize}
  \item \textbf{Transformation de variables}~: remplacer des
    variables qualitatives par les coordonn\'ees factorielles pour
    appliquer des m\'ethodes qui n\'ecessitent des variables
    quantitatives ($K$-means, r\'egression).
  \item \textbf{R\'eduction de dimension}~: travailler dans un
    espace de dimension r\'eduite pour att\'enuer le fl\'eau de la
    dimension (\emph{curse of dimensionality}).
  \item \textbf{Tandem clustering}~: appliquer une classification
    ($K$-m eans, CAH) sur les coordonn\'ees factorielles plut\^ot
    que sur les variables originales, ce qui am\'eliore g\'en\'eralement
    la qualit\'e de la partition.
  \item \textbf{Visualisation}~: les plans factoriels permettent de
    visualiser la partition obtenue et d'interpr\'eter les groupes.
\end{itemize}

\subsection{M\'ethodes multi-blocs}

Les m\'ethodes multi-blocs constituent un prolongement des analyses
factorielles classiques. Elles permettent d'analyser simultan\'ement
plusieurs tableaux de donn\'ees (blocs) mesur\'es sur les m\^emes
individus.

Les principales m\'ethodes sont~:
\begin{itemize}
  \item \textbf{DACP (Diagonalisation d'un Tableau de Covariance
    Par blocs)}~: g\'en\'eralisation de l'ACP \`a plusieurs blocs.
  \item \textbf{STATIS (Structuration des Tableaux \`a Trois Indices
    de la Statistique)}~: m\'ethode en quatre phases (interstructure,
    compromis, intrastructure, analyse du compromis).
  \item \textbf{AFM (Analyse Factorielle Multiple)}~: \'equivalent
    de l'AFDM pour les blocs de variables, avec une pond\'eration
    \'equilibrante des blocs.
\end{itemize}

\begin{important}
  Les m\'ethodes multi-blocs sont traitées en d\'etail dans une
  fiche d\'edi\'ee.
\end{important}

% ===================================================================
\section{Cartes de Kohonen (SOM)}
% ===================================================================

Les cartes autoorganisatrices de Kohonen (Self-Organizing Maps, SOM)
sont des r\'eseaux de neurones non supervis\'es qui r\'ealisent
simultan\'ement une r\'eduction de dimension et une classification.

\subsection{Principe}

Une carte de Kohonen projette les donn\'ees de haute dimension sur
une grille (g\'en\'eralement 2D) tout en pr\'eservant la topologie
des donn\'ees~: deux individus proches dans l'espace original
restent proches sur la grille.

\begin{itemize}
  \item \textbf{Grille}~: ensemble de neurones organis\'es selon
    une topologie (rectangulaire, hexagonale).
  \item \textbf{Vecteurs de r\'ef\'erence}~: chaque neurone $i$
    poss\`ede un vecteur $w_i$ de m\^eme dimension que les donn\'ees.
  \item \textbf{Apprentissage}~: on pr\'esente chaque observation
    $x$ et on d\'etermine le neurone vainqueur (celui dont le
    vecteur $w_i$ est le plus proche de $x$ au sens de la distance
    euclidienne). On met \`a jour $w_i$ et les neurones voisins
    dans la grille.
\end{itemize}

\subsection{Fonction de voisinage}

La mise \`a jour du vecteur du neurone vainqueur $c$ et de ses
voisins suit~:

\begin{equation}
w_i(t+1) = w_i(t) + \alpha(t) \, h_{ci}(t) \, (x(t) - w_i(t))
\end{equation}

o\`u~:
\begin{itemize}
  \item $\alpha(t)$ est le taux d'apprentissage (d\'ecroissant).
  \item $h_{ci}(t)$ est la fonction de voisinage (gaussienne
    ou ``chapeau mexicain''), qui d\'ecro\^it avec la distance
    entre les neurones $c$ et $i$ sur la grille.
\end{itemize}

\begin{remark}
  Les cartes de Kohonen sont trait\'ees en d\'etail dans une fiche
  d\'edi\'ee.
\end{remark}

% ===================================================================
\section{Conclusion}
% ===================================================================

Les m\'ethodes non-supervis\'ees forment un pilier essentiel de la
fouille de donn\'ees~:

\begin{itemize}
  \item La \textbf{classification} (partitionnement, hi\'erarchique,
    cartes de Kohonen) permet de d\'ecouvrir des groupes homog\`enes.
  \item Les \textbf{r\`egles d'association} identifient des liens
    fr\'equents entre items dans de grands volumes de donn\'ees.
  \item L'\textbf{analyse factorielle} (ACP, ACM, AFDM, multi-blocs)
    offre des visualisations synth\'etiques et sert de
    pr\'e-traitement pour la classification.
  \item La \textbf{classification de variables} compl\`ete la bo\^\i te
    \`a outils en permettant de regrouper des variables redondantes.
\end{itemize}

Ces m\'ethodes sont compl\'ementaires et peuvent \^etre combin\'ees
(s\'equence analyse factorielle $\rightarrow$ classification,
\emph{tandem clustering}) pour une exploration efficace des donn\'ees.

% ===================================================================
\section*{Questions cl\'es (QCM)}
% ===================================================================

\addcontentsline{toc}{section}{Questions cl\'es (QCM)}

\begin{enumerate}

  \item Quelle est la principale diff\'erence entre m\'ethodes
    supervis\'ees et non supervis\'ees~?
    \begin{enumerate}
      \item Les m\'ethodes non supervis\'ees n\'ecessitent une
        variable cible
      \item Les m\'ethodes supervis\'ees n\'ecessitent une
        variable cible
      \item Il n'y a pas de diff\'erence
      \item Les m\'ethodes non supervis\'ees sont plus rapides
    \end{enumerate}
    \textbf{R\'eponse~: b (les m\'ethodes supervis\'ees
    n\'ecessitent une variable cible).}

  \item Dans $K$-means, que minimise-t-on~?
    \begin{enumerate}
      \item L'inertie inter-classe
      \item L'inertie intra-classe
      \item La variance totale
      \item Le nombre de groupes
    \end{enumerate}
    \textbf{R\'eponse~: b (l'inertie intra-classe).}

  \item Quel crit\`ere de saut est le plus utilis\'e en CAH~?
    \begin{enumerate}
      \item Saut minimum
      \item Saut maximum
      \item Lien moyen
      \item Ward
    \end{enumerate}
    \textbf{R\'eponse~: d (Ward, qui minimise la perte d'inertie
    intra-classe).}

  \item Que mesure le support d'un itemset $X$~?
    \begin{enumerate}
      \item La probabilit\'e de $X$ sachant $Y$
      \item La fr\'equence de $X$ dans les transactions
      \item La force de l'association $X \Rightarrow Y$
      \item Le nombre d'items dans $X$
    \end{enumerate}
    \textbf{R\'eponse~: b (fr\'equence de $X$ dans les
    transactions).}

  \item Un lift $> 1$ indique~:
    \begin{enumerate}
      \item Une association n\'egative
      \item Une association positive
      \item Aucune association
      \item Une erreur de calcul
    \end{enumerate}
    \textbf{R\'eponse~: b (association positive~: les items
    apparaissent plus souvent ensemble que sous l'ind\'ependance).}

  \item Quelle analyse factorielle est adapt\'ee aux donn\'ees
    quantitatives~?
    \begin{enumerate}
      \item ACM
      \item ACP
      \item AFDM
      \item AFC
    \end{enumerate}
    \textbf{R\'eponse~: b (ACP -- Analyse en Composantes
    Principales).}

  \item Quelle analyse factorielle est adapt\'ee aux donn\'ees
    qualitatives~?
    \begin{enumerate}
      \item ACP
      \item ACM
      \item AFDM
      \item Toutes les r\'eponses
    \end{enumerate}
    \textbf{R\'eponse~: b (ACM -- Analyse des Correspondances
    Multiples).}

  \item Qu'est-ce que le \emph{tandem clustering}~?
    \begin{enumerate}
      \item Appliquer deux classifications successives
      \item Appliquer une analyse factorielle puis une
        classification
      \item Appliquer une classification puis une analyse
        factorielle
      \item M\'elanger deux algorithmes de classification
    \end{enumerate}
    \textbf{R\'eponse~: b (analyse factorielle puis
    classification).}

  \item Quelle propri\'et\'e permet l'\'elagage dans l'algorithme
    Apriori~?
    \begin{enumerate}
      \item La monotonie du support
      \item L'anti-monotonie du support
      \item La sym\'etrie de la confiance
      \item La transitivit\'e du lift
    \end{enumerate}
    \textbf{R\'eponse~: b (anti-monotonie~: tout sur-ensemble
    d'un itemset non fr\'equent est non fr\'equent).}

  \item Dans les cartes de Kohonen, que pr\'eserve la projection~?
    \begin{enumerate}
      \item Les distances exactes entre individus
      \item La topologie des donn\'ees
      \item La variance des variables
      \item L'ordre des individus
    \end{enumerate}
    \textbf{R\'eponse~: b (la topologie~: deux individus proches
    dans l'espace original restent proches sur la grille).}

\end{enumerate}

% ===================================================================
\begin{thebibliography}{9}
% ===================================================================

\bibitem{AudigierNiang2026}
  V.~Audigier, N.~Niang.
  \newblock \emph{STA211~: M\'ethodes non-supervis\'ees
    (introduction)}.
  \newblock Cours CNAM, 20 f\'evrier 2026.

\bibitem{MacQueen1967}
  J.~MacQueen.
  \newblock \emph{Some methods for classification and analysis of
    multivariate observations}.
  \newblock Proc. 5th Berkeley Symp. Math. Statist. Probab., 1967.

\bibitem{AgrawalSrikant1994}
  R.~Agrawal, R.~Srikant.
  \newblock \emph{Fast algorithms for mining association rules}.
  \newblock VLDB, 1994.

\bibitem{Kohonen1982}
  T.~Kohonen.
  \newblock \emph{Self-organized formation of topologically correct
    feature maps}.
  \newblock Biological Cybernetics, 43(1):59--69, 1982.

\bibitem{Rakotomalala}
  R.~Rakotomalala.
  \newblock \emph{Classification ascendante hi\'erarchique}.
  \newblock \url{https://eric.univ-lyon2.fr/ricco/cours/cours/CAH.pdf}

\end{thebibliography}

\end{document}
